\chapter{Annexes}
\markboth{ANNEXES}{}

\section{Guide d'installation et d'exécution}
Le système requiert Python~3.10 ou supérieur.

\begin{lstlisting}[caption={Installation et lancement.}]
# 1. Environnement virtuel
python -m venv .venv
source .venv/bin/activate        # macOS / Linux

# 2. Dependances
pip install -r requirements.txt

# 3. Configuration (fournisseur de LLM)
cp .env.example .env             # renseigner GROQ_API_KEY, ou LLM_PROVIDER=ollama

# 4. Lancement de l'interface web
uvicorn src.ui.app:app --reload
# -> ouvrir http://127.0.0.1:8000
\end{lstlisting}

Il est aussi possible d'exécuter le pipeline en ligne de commande sur un
fichier, ou de générer le jeu de données d'exemple :

\begin{lstlisting}[caption={Exécutions alternatives.}]
# Pipeline en ligne de commande
python -m src.pipeline data/sample_sales.csv

# Generation d'un jeu de donnees d'exemple
python scripts/generate_sample_data.py --out data/sample_sales.csv
\end{lstlisting}

\section{Extrait : matrice de permissions (\texttt{config.yaml})}
La configuration ci-dessous illustre le contrôle d'accès par rôle appliqué
par le serveur MCP.

\begin{lstlisting}[caption={Définition d'un agent et de ses outils autorisés.}]
agents:
  kpi:
    description: >
      Metriques et insights. Curent et calculent les KPIs, puis
      detectent anomalies et comparaisons de segments.
    tools:
      - profile_dataset
      - suggest_kpis
      - compute_kpis
      - detect_anomalies
      - compare_segments
\end{lstlisting}

\section{Extrait : boucle d'appel d'outils d'un agent}
\begin{lstlisting}[language=Python,caption={Cœur de la boucle agentique (\texttt{agents/base.py}).}]
for step in range(1, self.max_steps + 1):
    response = self.llm.chat(messages, tools=tools, temperature=self.temperature)
    if not response.tool_calls:
        # relance avec incitation si l'agent n'a encore appele aucun outil
        if all_tool_calls or not tools or nudges_used >= self.nudge_budget:
            final_text = response.content; break
        messages.append(nudge_message); nudges_used += 1; continue
    for tc in response.tool_calls:           # execution des appels d'outils
        result = self.server.call_tool(store, self.role, tc.name, tc.arguments)
        messages.append(tool_result_message(tc, result))
\end{lstlisting}

\section{Arborescence du projet}
\begin{lstlisting}[caption={Organisation des sources.}]
src/
  core/        # moteur deterministe : profiler, kpi_engine,
               # dashboard_builder, insights, validation
  llm/         # client LLM (Groq / Ollama)
  mcp_server/  # serveur MCP : outils + permissions
  agents/      # les 7 agents specialises
  ui/          # interface FastAPI + gabarits + rendu du rapport
  pipeline.py  # orchestrateur de bout en bout
tests/         # 22 tests automatises
config.yaml    # matrice de permissions + seuils du moteur
\end{lstlisting}
